\section{Architecture}\label{sec:4}

\emph{This section carries three subsections beyond the template's four, so its numbering runs
one ahead from 4.3 onwards. The template's \textbf{4.3 Component diagram} is \S\ref{sec:4.4} here, and its
\textbf{4.4 Design decisions and alternatives} is \S\ref{sec:4.7}. Nothing is omitted: \S\ref{sec:4.3} adds the
authority separation, \S\ref{sec:4.5} the two code-level diagrams, and \S\ref{sec:4.8} how the architecture grows,
all three of which the architecture needs and none of which the template has a slot for.}

\subsection{Context diagram (C4 Level 1)}\label{sec:4.1}

The context view fixes the \textbf{system boundary}. The system reads and commands the building
only through BuildSim; the technician receives a work order and, by repairing the damper, is the
only actor who can clear a latched case; occupants are sensed, not connected. Nothing internal
appears.

\begin{center}\includegraphics[width=0.58\textwidth]{figures/context.pdf}\end{center}
\figcap{C4 System Context (Level 1), from \texttt{project/diagrams/context.d2}. Person and
software-system elements; every arrow is a directed relationship labelled with its purpose and
protocol.}

\subsection{Container diagram (C4 Level 2)}\label{sec:4.2}

Each box is one Docker container, one service in \texttt{docker-compose.yml} and one folder under
\texttt{project/}. That 1:1 mapping is a hard constraint of the implementation plan, not a
coincidence, and it is what keeps the diagram from drifting from the code: adding a box means
adding a folder and a service, which is why the fan-speed command is applied by
\texttt{damper-actuator} rather than by a container of its own (Section \ref{sec:4.7}).

\textbf{\factComposeServicesCap{} compose services, \factAuthoredServices{} of them project-authored}: \texttt{physics-simulator},
\texttt{co2-sensor}, \texttt{damper-actuator}, \texttt{dcv-controller}, \texttt{fault-monitor},
\texttt{supervisor} and \texttt{dashboard}, plus \texttt{buildsim} from the \texttt{buildingsim/}
folder of the same repository and the third-party images \texttt{mosquitto} and
\texttt{timescaledb}, configured under \texttt{project/deploy/}. Two folders are deliberately
\emph{not} containers (\texttt{shared/} is a library and \texttt{evaluation/} the offline harness),
which is the entire reason the folder set is larger than the service set, and the mapping stays
1:1 for all \factAuthoredServices{}.

\begin{center}\includegraphics[width=\textwidth,height=0.70\textheight,keepaspectratio]{figures/container.pdf}\end{center}
\figcap{C4 Container diagram (Level 2), from \texttt{project/diagrams/container.d2}. The dashed
boundary is the system in scope; BuildSim is external; every arrow carries its protocol. The
figure is service-for-service identical to \texttt{project/docker-compose.yml}.}

The loop the figure describes is the one that makes this a cyber-physical system rather than a
dashboard. \texttt{physics-simulator} reads actuator state from BuildSim, advances the physics and
writes sensor values back; \texttt{co2-sensor} reads those values and publishes them;
\texttt{dcv-controller} commands a damper position; \texttt{damper-actuator} applies it to
BuildSim; and that actuator change alters the environment the simulator advances on its next tick.
The outer loop rides the same bus one order of magnitude slower.

The technician writes into the system in two places. One reports the equipment
serviced to \texttt{dashboard}, which is the only exit from a latched case; the other is a small
HTTP fault control surface on \texttt{physics-simulator} (\texttt{POST}, \texttt{GET} and
\texttt{DELETE} on \texttt{/faults}) that injects, lists and clears a fault in the running
plant. The second is what makes a recovery demonstrable rather than merely described: the
scenario can be changed and the damper repaired without restarting the stack. Both ports are
published on \texttt{127.0.0.1} only.

\subsection{Authority separation: who may act}\label{sec:4.3}

\textbf{\texttt{fault-monitor} raises events and never commands. \texttt{supervisor} is the only
component that acts.} This is the single most important structural decision in the system and it
is worth being explicit about why, because the two halves could obviously have been one process.

\begin{itemize}
\tightlist
\item \textbf{A crashing monitor must not be able to stop ventilation.} Diagnosis is the
  speculative, model-dependent, newest and least-trusted code in the system. Control is the part
  that keeps people breathing. If they share a process, a null pointer in the classifier takes the
  damper commands down with it. Split, the worst a dead monitor can do is stop producing
  diagnoses, and ventilation continues on the inner rule.
\item \textbf{A wrong diagnosis must not be able to command anything by itself.} The monitor's
  output is a belief about the world, derived from a model I wrote about physics I also wrote. A
  belief that can actuate is a belief that can degrade a healthy room. Separating them means every
  action passes through a component whose only input is a \emph{confirmed} fault, and that
  component is small enough to read in one sitting.
\end{itemize}

The separation is enforced structurally rather than by convention: the monitor's container has no
publisher for any \texttt{commands/*} topic, and the single exception (the per-room active probe,
which is a damper command) is owned by the supervisor precisely so that the rule has no
exception. The same reasoning fixes ownership of the fan speed: \texttt{commands/\{ahu\}/fan-speed}
is the compensation lever, so the supervisor owns it and the inner loop never writes it.

\begin{longtable}[]{@{}
  >{\raggedright\arraybackslash}p{(\columnwidth - 6\tabcolsep) * \real{0.170}}
  >{\raggedright\arraybackslash}p{(\columnwidth - 6\tabcolsep) * \real{0.310}}
  >{\raggedright\arraybackslash}p{(\columnwidth - 6\tabcolsep) * \real{0.135}}
  >{\raggedright\arraybackslash}p{(\columnwidth - 6\tabcolsep) * \real{0.385}}@{}}
\toprule
\textbf{Container} & \textbf{Responsibility (one job)} & \textbf{Tech} & \textbf{Why it is a
separate process} \\
\midrule
\endhead
\texttt{physics-simulator} & Advance the CO₂ balance, the occupancy schedules, the air handler and
the injected fault; write sensor values to BuildSim; take a fault in or out at runtime & plain Java 25 & It is the environment, not
the system. Nothing else may write a sensor value, or the loop would not close through the
building \\
\texttt{co2-sensor} & Read six CO₂ values and the unit's fan power from BuildSim; judge data
quality; publish & plain Java 25 & Emulates the devices. Restartable without touching control;
its failure must produce silence, which is what R-4 detects \\
\texttt{damper-actuator} & Validate and apply HVAC commands to BuildSim: six dampers and the one
fan speed & plain Java 25 & The only writer of actuator state. A separate failure domain from
whoever decided the command \\
\texttt{dcv-controller} & Inner loop: damper position from CO₂, with hysteresis and a freshness
check & Quarkus JVM & Must survive the monitor and the supervisor dying. Its rule changes on a
different schedule from the diagnosis \\
\texttt{fault-monitor} & Observer, residual, one-sided CUSUM, state machine, classifier.
\textbf{Events only} & Quarkus JVM & See above: the least-trusted code must not hold actuator
authority \\
\texttt{supervisor} & Act on a confirmed fault: degrade, compensate, work order & Quarkus JVM &
The only component that commands. Small, auditable, and its input is already a decision \\
\texttt{dashboard} & Residual and fault state as BuildSim room layers, plus its own page and JSON
& Quarkus JVM & Presentation must not be able to slow control. Every dependency is optional so the
page still serves when the rest is down \\
\texttt{mosquitto} & Decouple producers from consumers & Mosquitto 2 & One source feeds many
consumers; a dead publisher goes silent rather than blocking \\
\texttt{timescaledb} & Hot store for readings, commands, residuals, decisions, strategies &
TimescaleDB pg16 & The residual window query is the monitor's hot path and the journal is R-5's
evidence \\*
\texttt{buildsim} & Building state and the 3D view & Go (unmodified) & External system; pinned by
being in the same repository (ADR-0007) \\
\bottomrule
\end{longtable}

\subsection{Component diagram (C4 Level 3) and the code structure}\label{sec:4.4}

\texttt{fault-monitor} is the one container complex enough to deserve its own picture, and
\textbf{Appendix \ref{sec:B.1}} draws it: every box one \emph{component}, labelled with the exact type name so
a reader can grep for it, with the \texttt{boundary} group holding every framework and transport type
and \texttt{control} and \texttt{entity} depending on nothing outward. What matters in the body is
the rule the figure obeys.


\textbf{BCE inside hexagonal.} Packages are organised by \textbf{business component}
(\texttt{ventilation/}, \texttt{diagnosis/}, \texttt{supervision/}, \texttt{physics/},
\texttt{devices/}, \texttt{dashboard/}), and each business component has the same three layers:

\begin{itemize}
\tightlist
\item \texttt{entity/}: the domain model as Java records, with invariants enforced in compact
  constructors. \texttt{FanSpeed} cannot hold 1.4; \texttt{DamperPosition} cannot hold $-0.1$;
  \texttt{AirHandlingUnit} cannot be commanded outside its range. An illegal state is
  unconstructible rather than merely untested.
\item \texttt{control/}: the business logic: the observer, the residual accumulator, the state
  machine, the classifier, the compensation solver, the hysteresis rule. No framework type appears
  here, which is why all of it is unit-testable without a broker, a database or a container.
\item \texttt{boundary/}: every adapter: Paho for MQTT, JDBC with native SQL for TimescaleDB, an
  HTTP client for BuildSim, Quarkus scheduling and health.
\item \texttt{event/}: \textbf{a fourth package, in \texttt{shared} only}, and it is named here
  because a reader who opens \texttt{shared/src/main/java/se/ltu/d7065e/shared/} finds three
  directories and this description otherwise promises two. It holds \texttt{DomainEvent}, its eight
  permitted records and \texttt{Topics}. It is not a boundary (it carries no framework import),
  and it is not one component's entity layer, because the alternative to a shared package is a
  dependency from every subscriber onto every publisher. ADR-0005 and ADR-0006 record the choice.
\end{itemize}

\textbf{The dependency rule} is what makes this hexagonal rather than merely three-layered: the
arrows point \emph{inward}. \texttt{control} and \texttt{entity} never reference
\texttt{boundary}; instead the core declares \textbf{ports} in domain vocabulary. There are seven,
all of them in \texttt{shared}'s \texttt{control} package and nothing else in it:
\texttt{ReadingStreamPort}, \texttt{EventPublisherPort}, \texttt{ResidualJournalPort},
\texttt{BuildingStatePort}, \texttt{BuildingOverlayPort}, \texttt{OccupancyPort} and
\texttt{ClockPort}. And \texttt{boundary} implements them. One adapter may implement more than
one: \texttt{MqttAdapter} is a single class satisfying both \texttt{ReadingStreamPort} and
\texttt{EventPublisherPort}, because subscribing and publishing are one broker session. The payoff is concrete and is not architectural tidiness for
its own sake: because the clock is a port, the evaluation harness can drive seven simulated days
through the same classes the containers run, in minutes; and because the reading stream is a port,
every control-layer test runs with an in-memory fake instead of Docker.

Two honest notes on the structure. There is \textbf{no DTO or mapper layer}: domain records are
serialised directly, and translation exists only where the schema is not mine (BuildSim and the
journal). That is justified by lock-step single-author deployment and is wrong the moment any
consumer deploys independently: it makes a domain rename a wire-format break. And
\texttt{shared/} holds both the domain language and the three infrastructure adapters, so it is a
library two kinds of thing live in; splitting it would add a module for no behavioural gain today.

\subsection{Code diagrams (C4 Level 4), and why only two}\label{sec:4.5}

Two places in this system are claims about the \emph{shape of the code} rather than about the
physics, and at Level 3 both are hidden inside a single box. The first is the diagnosis chain inside
\texttt{fault-monitor}: that \texttt{runCycle()} has two phases and the order is load-bearing, that
the published residual is not the accumulated one, that \texttt{serviced()} now has a transport, and
that \texttt{DAMPER\_STUCK\_OPEN} exists as a classification without being detectable. The second is
\texttt{CompensationSolver.solve}, whose claim is that holding five healthy flows while multiplying
the sixth is a \textbf{closed form rather than a search}: straight-line code with no convergence
test. C4's own guidance is that the code level is optional and worth drawing only where it earns its
page, so both figures and the four claims they carry are in \textbf{Appendix \ref{sec:B.2}} rather than here.
Everything else in the system is either a plain adapter, where the port name is already the whole
story, or logic whose interesting content is arithmetic that Sections \ref{sec:6} and \ref{sec:7} give in full.

\subsection{The documentation approach: C4 for the views, arc42 for the frame}\label{sec:4.6}

Two frameworks are used, for two different jobs, and each is weak where the other is strong.
\textbf{C4 supplies the structural views}: a numbered zoom from context to containers to components
to code, with a fixed element vocabulary at each level, so a reader can tell what kind of thing a box
is without reading its label twice. \textbf{arc42 supplies the document frame}: twelve numbered
sections that force the questions a structural diagram cannot answer (constraints, cross-cutting
concepts, quality scenarios, risks and technical debt). C4 has no place to put a measured quality
scenario or an honest debt register, and arc42 names a ``building block view'' without saying what a
building block may be.

The arc42 document is \texttt{project/docs/arc42.md}, and it deliberately does \emph{not} repeat
this report: where the report says something, arc42 cites the section and summarises it, and its
value is the two things this report has no natural home for: the quality tree with its measurable
scenarios (arc42 §\ref{sec:10}) and the debt register as of a stated date (arc42 §\ref{sec:11}). The C4 diagrams are
shared rather than duplicated: both documents reference the same nine files in
\texttt{project/diagrams/}.

\subsection{Design decisions and alternatives}\label{sec:4.7}

Each row states the choice, the alternative genuinely considered, the requirement it serves and
the price paid. A decision with no rejected alternative is only an assumption. The full records
are in \texttt{project/decisions/}.

\begin{longtable}[]{@{}
  >{\raggedright\arraybackslash}p{(\columnwidth - 8\tabcolsep) * \real{0.165}}
  >{\raggedright\arraybackslash}p{(\columnwidth - 8\tabcolsep) * \real{0.155}}
  >{\raggedright\arraybackslash}p{(\columnwidth - 8\tabcolsep) * \real{0.185}}
  >{\raggedright\arraybackslash}p{(\columnwidth - 8\tabcolsep) * \real{0.285}}
  >{\raggedright\arraybackslash}p{(\columnwidth - 8\tabcolsep) * \real{0.210}}@{}}
\toprule
\textbf{Decision} & \textbf{Chosen} & \textbf{Alternatives considered} & \textbf{Justification
(requirement)} & \textbf{Trade-off accepted} \\
\midrule
\endhead
Language and runtime (ADR-0002) & Java 25 on Quarkus, one container per component & \textbf{Go,
which the course recommends}; Erlang/OTP & The course allows another language \emph{if strongly
motivated}, and the motivation is the assessment itself: the oral is individual, unaided and
caps the grade at what can be defended at a whiteboard, so fluency is an architectural input
rather than a preference. Erlang was weighed properly and fails on the requirement: OTP
supervision runs \emph{inside} the BEAM, so using it as intended puts N sensors in one Unix
process and violates the one-process-per-device rule, while one node per sensor is the heaviest
possible host for a single actor & A larger baseline container footprint than Go on the trivial
sensor processes. ADR-0002 undertook to \textbf{measure} this rather than argue it, because
footprint per container feeds the breaking-point experiment directly (§\ref{sec:11}, R-6); the figure was
never taken, so it stands as accepted rather than measured \\
Coupling topology (ADR-0008) & One shared air handler, a damper per room, proportional split &
Six independent rooms with a fan each; shared handler \emph{plus} inter-room air exchange &
Creates a real cross-room signal (plant fan power) and makes R-7 something the supervisor must
earn rather than something the topology gives away (R-2, R-7) & A seventh simulated element and
its invented fan curve; compensation cannot reach a damper stuck fully shut \\
Monitor/supervisor split & Two containers, events one way & One ``autonomous service'' &
A crashing monitor cannot stop ventilation; a wrong diagnosis cannot command (R-4). \textbf{Measured
rather than assumed}: T-27 kills each component in turn; 23 of 23 assertions in one JVM, and at
container level five of six rooms reach \texttt{BLIND} within 12 s while ventilation continues
(§\ref{sec:11.2}, Appendix \ref{sec:G}) & One more process, and a topic hop in the latency budget \\
Detection statistic & One-sided CUSUM on the residual, \texttt{k} = 30 ppm, \texttt{h} = 150 ppm
& Raw residual threshold; two-sided CUSUM; learned one-class model & Accumulation catches a small
persistent bias long before a threshold fires, and the spike measured 16.7 σ separation (R-1, R-3).
\textbf{The learned alternative was not merely reasoned about, it was built and measured}: on a
held-out fault-free week T-29's contender raises \textbf{11.25} false alarms per week against the
rule's \textbf{0.55}, and it was dropped on that evidence (§\ref{sec:11.2}, Appendix \ref{sec:F}) & One-sided means a
sensor drifting \emph{low} is invisible (§\ref{sec:13}) \\
Inner control law & Plain rule with hysteresis and a freshness check & PID; model-predictive
control & The outer loop must model the inner one, so inner complexity becomes observer model
error (R-1) & Less efficient ventilation than an optimiser would give \\
Clearing policy & Latch; only an explicit service event clears & Clear when the residual falls &
Compensation shrinks the residual that found the fault, so a self-clearing monitor would
oscillate with a broken damper still broken (R-2) & A false positive is sticky, and
a false positive costs one technician visit and one HTTP request (§\ref{sec:13}) \\
Fan-speed ownership & \texttt{supervisor} publishes it; \texttt{damper-actuator} applies it & A
\texttt{fan-actuator} container of its own & The fan is the compensation lever, so the inner loop must
not write it; a container of its own for one PUT is not worth a box, a folder and a service (plan
constraint 4) & \texttt{damper-actuator}'s name is now narrower than its job \\
Transport & MQTT pub/sub for readings, commands and fault events & Direct REST between components;
Kafka & One sensor feeds controller, monitor, store and dashboard; a dead publisher goes silent
rather than blocking (R-4) & A broker to run; and without retained replay a monitor restart loses
its residual history \\
Service state across a restart (ADR-0012) & One JSON document per service on its own volume &
Replaying the journal to rebuild state; a second database (MongoDB was adopted and then rejected);
accepting the loss & The journal is an append-only record of \emph{what happened} and this is
\emph{what is believed now}: different lifetimes, and one of the three things needed (the pinned
damper belief) exists nowhere in the journal but as prose. A restarted monitor kept no latch, and a
restarted supervisor \emph{erased} live work orders & Host-local: the file follows the volume, not
the service, so a multi-host deployment must replace it. No history and no query: correct for
state, wrong for evidence, which is why the journal is untouched \\
Hot store & TimescaleDB with native SQL over six hypertables & Plain PostgreSQL; files; an ORM &
The residual window query per \texttt{caseId} is the hot path, and R-5 needs provenance on every
row & Native SQL is hand-maintained; one more stateful service \\
Runtime & Java 25; Quarkus JVM for the four scheduled services, plain \texttt{java -jar} for the
three device processes & Quarkus everywhere; Spring Boot; native images in development & Scheduling,
config and health are worth a framework; a 40-line sensor loop is not (ADR-0002) & Two startup
styles to understand in one repository \\
Architecture style & BCE inside hexagonal, ports in the core & Layered by technology; framework
annotations in the core & Lets the evaluation harness drive the same classes through a simulated
clock, and keeps control-layer tests container-free (R-6, R-3) & More indirection than a small
system strictly needs \\
Fault injection point & Between the command and its physical effect & At the sensor; at the
controller & The command record and the physics must disagree, or there is nothing to detect
(R-1) & The fault model is mine, and so is the detector (§\ref{sec:13}) \\
Healthy-room airflow bound (ADR-0009) & One symmetric 5 \% allowance, kept, with the two sides
separated in the code because they are different quantities & An asymmetric over-supply allowance
at 0.10, 0.20, 0.43 or unbounded; a per-room trim instead of one common factor & Swept closed-loop,
the allowance buys R-2 zero occupied minutes at every setting (0.024 and 0 of 5 seeds throughout)
because with every healthy damper on its floor the fan lever is not selective and the plant caps it
at 1.43 against the 2.0–3.3 the regime asks (R-2, R-7) & R-2 keeps a named failure instead of a
bought pass, at 134 ppm of mean concentration the plant could have delivered \\
Plausibility guards are bounds, not filters (ADR-0010) & Drive the guarded lower arm off the
outdoor clamp, and re-anchor the other five rooms when a confirmation pins a damper & Gate both
arms on the judging window; never let a re-anchor lower the ceiling; clamp at the design generation
alone; raise $h$; fix the lower arm only & \texttt{correctTowards} pulled both guards 30 \% of the
way onto every reading, so each was a filter and each arm ran on a filter error rather than on a
residual; correcting the lower arm's half halves R-3's rate and moves no other requirement
(R-3) & R-3 is still not met, and the variant that does meet it is left on a branch because it
costs R-1's closed-loop arm 15 of 20 against 1 of 20, T-27 and R-7 \\
\bottomrule
\end{longtable}

\subsection{How this architecture grows}\label{sec:4.8}

Three growth directions are worth separating, because the architecture answers them very
differently: \textbf{more zones on one plant}, \textbf{more plants or more buildings}, and
\textbf{new or replaced services}. For the first two the limits are measured and the bottleneck is
named. For the third the honest answer is not a list of seams but the one boundary that
\emph{looks} like a seam and is not, so this subsection ends there rather than beginning there.

\textbf{More zones on one plant.} Two ceilings exist and they are never quoted as one. The
harness reaches \textbf{25\,600 zones at a p95 of about 6 s} inside the 60 s interval, fits
p95 = $1.196\times10^{-8}\,n^{1.97}$ with $r^2$ = 0.9991, and extrapolates the 60 s crossing to
roughly \textbf{86\,000 zones} (§\ref{sec:11}, R-6). The \emph{deployed} ceiling is \textbf{58 zones}, and it
is BuildSim's floorplan (the simulator will not host equipment in a room its level-0 plan does not
have), not anything about the monitor. The exponent 1.97 is the useful part: the saturated element
is \texttt{AhuFlowModel.believedFlowTo}, which sums every damper's opening to find the denominator
of the proportional split and is therefore $\Theta(n)$ per zone and $\Theta(n^2)$ per cycle. The
harness names it in its own output rather than leaving it to be inferred
(\texttt{ScalingExperiment.java:122}). It is also the cheapest thing in the system to fix: the
denominator is one scalar that changes only when a damper moves, so caching it and invalidating on
change makes the cycle $\Theta(n)$ and moves the crossing out by orders of magnitude. That work is
not done because 58 zones is the ceiling that binds, and optimising past a limit the deployment
cannot reach would be sophistication with nothing behind it.

\textbf{More plants, or more buildings.} \textbf{The shard boundary is one air-handling unit}, and
it is a shard rather than a partition for a structural reason: the only cross-room term in the
whole design is the flow-split denominator above, and that sum runs over one plant's dampers.
Two air handlers share no term, so two monitors over two AHUs compute exactly what one monitor over
each would: no coordination, no distributed state, no consensus. The wire is already shaped for
it: \texttt{Topics.fanPower(AhuId)} and \texttt{Topics.fanSpeedCommand(AhuId)} take the handler as a
parameter, so \texttt{readings/\{ahu\}/fan-power} is per-plant today and a second AHU needs no
topic design at all. One \texttt{fault-monitor}, one \texttt{supervisor} and one
\texttt{dcv-controller} per AHU, sharing the broker and the journal, is the whole scaling story for
a building, and the per-zone tables already carry the zone, so the journal needs no change either.
What is \emph{not} shaped for it is named below.

\textbf{New or replaced services.} Every boundary at which an implementation can be swapped is
inventoried in \texttt{project/diagrams/seams.md}, which separates two kinds that are routinely
confused. A \textbf{contract seam} is run-time and language-agnostic (an MQTT topic plus the
payload record on it), so anything publishing the right topic with the right payload is
substitutable whatever language it is written in, and the far side cannot tell; this is where a
physical sensor or a real DDC panel is installed. A \textbf{port seam} is compile-time and
same-JVM: a Java interface in \texttt{shared/control}, declared in the core's vocabulary rather
than mirroring an adapter's API (ADR-0005), swapped in a composition root with no change to
\texttt{control} or \texttt{entity}. \textbf{A port seam cannot install a device, and a contract
seam cannot be unit-tested with a stub}, which is why the system has both.

\begin{center}\includegraphics[width=0.92\textwidth]{figures/seams-overview.pdf}\end{center}
\figcap{The contract seams, from \texttt{project/diagrams/seams-overview.d2}. Hexagons are
contract seams at the wire; the tier colouring is in the legend below. The per-boundary figures
(\texttt{seams-plant}, \texttt{seams-transport}, \texttt{seams-storage}, \texttt{seams-clock},
\texttt{seams-topology}) carry the port seams and the non-seams.}

\begin{center}\includegraphics[width=0.52\textwidth]{figures/seams-legend.pdf}\end{center}
\figcap{Tier 1 clean: replace one side, nothing else changes. Tier 2: genuinely substitutable, but
something named must change with it or be inherited. Tier 3: looks like a boundary on the container
diagram, and crossing it would mean changing code on both sides. \textbf{The Tier 2 conditions and
the Tier 3 entries are the content}. Twelve seams all claimed clean would be a prettier set of
figures and a worse one.}

Two questions an examiner is entitled to ask are each answerable by pointing at one line.
\emph{Where do we cut to replace the virtual CO\textsubscript{2} sensor with a real one in the
ceiling?} At \textbf{K1}, the wire, where \texttt{readings/\{zone\}/co2} carries
\texttt{ReadingObserved}, not at \texttt{BuildingStatePort}, because that is a Java interface on
this classpath and a device in a ceiling is not a Java class. K1 is Tier 2 and the replacement
inherits three stated conditions: it must judge its own data quality, because nothing downstream
can; it must publish fan power too, or nothing does; and it must stamp \texttt{ObservedAt} on the
shared timeline. \emph{Where do we cut to replace the supervisor?} At \textbf{K2}, whose three
published topics are the complete \emph{control} contract, provided the replacement also inherits
two obligations that are not topics at all: it must be the only writer of BuildSim's alert
collection, and it must journal its own decisions.

\textbf{And the part that is not a seam, which is the real answer to how this grows.}
\texttt{Topology.java} hard-codes the plant: the six rooms with their areas and design flows, the
nominal flow and power, the normal fan speed, the damper floor, the outdoor and comfort
concentrations, the generation per person, both noise sigmas, and the CUSUM slack and limit. It is
an interface of static members in \texttt{shared/}, \textbf{with no port in front of it and no
injection point}, and \textbf{34 main-source files import it}. And the readers are
\texttt{control} classes rather than adapters, which is precisely what makes it unsubstitutable.
A real corridor has measured design flows, a measured air handler and a measured sensor noise
floor, so every one of those numbers is edited in one file in the library all \factAuthoredServices{}
services link, which recompiles and redeploys all \factAuthoredServices{}. \texttt{Topology.AIR\_HANDLER} is likewise a single
literal, \texttt{AhuId.of("AHU-A120")}, so the per-AHU sharding the topics already support is
blocked by one constant rather than by any part of the design.

The sharpest consequence is not the redeploy. \textbf{The CUSUM slack and limit are tuned to the
sensor noise sigma} and the file says so in place: a real sensor with a different noise floor does
not merely change a constant, it \textbf{invalidates the spike that measured $k$ and $h$ together}
(16.7 σ separation, 0 of 20 false alarms, §\ref{sec:11.1}) until that spike is redone against the real
instrument. That is the honest cost of taking this architecture to a real building, and it is a
\emph{commissioning} cost rather than a software one. If one thing were done first it would be
this: put the plant behind a configuration port so the numbers are data rather than code, and make
the tuning of $k$ and $h$ an explicit function of the measured noise floor instead of a pair of
constants that happen to match it.\label{sec:4end}
